% s-292-p292-m001; original PDF 292, 293, 294, 295
\Aphorism{s-292-p292-m001}{353}
Es ist ein Zustand auf Erden denkbar, da die „Wahrheit“ sich biologisch so stark festgesetzt und einverleibt haben wird, dass eine abweichende Meinung zu den pathologischen Fällen gehören würde. In diesem „Gottesreich“ werden die Ansichten, Glaubenssätze, Resultate für alle Menschen übereinstimmen, und nur in der Methodik, in der Bildungs- und Entwicklungsgeschichte dieser Resultate wird noch etwas wie Persönlichkeit und individuelle Eigenart möglich sein. Schon jetzt haben wir einen Vorgeschmack solcher Zustände; schon jetzt uniformirt man gründlich genug in allen Meinungen und überlässt dem persönlichen Belieben nur die Vorgeschichte, die Genealogie der Meinungen — ähnlich wie man dem Adel zwar den privaten Luxus seiner Stammbäume gönnt, ihn aber \Foreign{de jure} in den bürgerlichen Gesammtverband einreiht. So dürfen wir auch heute zwar noch unsere eigenen Curven durchlaufen, aber die Endpunkte sind grossentheils fixirt: „Alle auf allen Wegen gelangen zu mir“ — kann die „Wahrheit“ mit gutem Gewissen sagen. So kommt es, dass Einer, was den blossen Bestand an fertigen, handgreiflichen Meinungen anbelangt, mit Vielen verwechselt werden kann, von denen er sich innerlich so abgetrennt wie nur möglich fühlt. Nehmen wir einen Fall: den \Sperr{Optimismus}. Dass ein Philosoph wie Nietzsche zum Leben bejahend steht, unterscheidet ihn äusserlich nicht einmal vom ersten besten Rechtfertigungsvirtuosen aus der Hegelschen Vernünftigkeitsschule oder vom vergnügten Philister, der seine zufällige Behaglichkeit zum \Greek{πάντα} \Greek{καλὰ} \Greek{λίαν} verallgemeinert. Man kann dem Dasein décharge ertheilen, weil man seine Bilanz nicht nachzurechnen versteht oder nur sub specie sui nachrechnet; man kann über den Werth des Lebens urtheilen wie der Eingeweidewurm über den Werth des Menschen urtheilen würde. „Ein durchaus wohlschmeckendes Thier“, sagte ein alter Eingeweidewurm, „das nichts dafür kann, wenn unsere pessimistischen Feinschmecker seiner überdrüssig werden. Für vernünftige Ansprüche ist dennoch der Unterleib des Menschen die beste aller möglichen Welten.“ Dem Philosophen wird die Selbstgefälligkeit, mit der das kleinste Ungeziefer aus seinen entozoologischen Erfahrungen und Werthschätzungen heraus den Advocaten des Lebens macht, immer etwas komisch bleiben; dergleichen Urtheile über das Leben, die eine Stellung ausserhalb des Lebens voraussetzen, sollten billigerweise nur einem \Sperr{Ektozoon} zustehen. Aber, wie gesagt, rein als Meinung und Theorem betrachtet, wäre die Weltbejahung des freiblickenden Philosophen mit der analogen Gefühlsäusserung der blinden Schmarotzer und Eingeweidewürmer ebenso leicht zu verwechseln wie mit dem arglosen Vernünftigkeitsglauben der Winkelmänner und Eckensteher, die nicht um ihre Ecke sehen: der Unterschied liegt allein in der \Sperr{Vergangenheit} beider Optimismen!
\par

Der eine war nie etwas anderes als Optimismus, unbefangener, kritikloser, blauäugiger und blondlockiger Optimismus, sei es als „ideologischer Überbau“ zum engen Wohlbehagen des petit bourgeois oder als aurea mediocritas treuherziger Menschlein, die am Leben nichts Fragwürdiges finden können. Der andere, aus den Wildnissen und Höhlengängen einer pessimistischen Jugend hervorgekrochen, ein Optimismus \Foreign{malgré soi}, ein Jasagen durch doppelte Verneinung, eine Liebe zum Leben auf Grund — eines eisigen Verdachts gegen das Leben, schaudernder Einblicke in den Unsinn des Lebens, ausschweifender Paroxysmen des Mitleids mit allem Lebenden! Eine harte und gewaltsame Zwangsperspective, eine peinliche Frage, bei der eigentlich der Richter der Gefolterte ist, ein grundsätzliches „Zurückgehen auf die Quellen“, auf die letzten Voraussetzungen, auf die hin das Leben zu Recht besteht, ein höchster Grad von Unabhängigkeit gegenüber den kategorischen Imperativen des Lebens — und als Ende dieser langen, langwierigen Kritik der „Lebenslüge“ eine neue, skeptische, beinahe wissenschaftliche Liebe zum Leben: der pessimistische Sturm hat die Luft gereinigt und die Wolken von den nahen Berggipfeln hinweggefegt! Wieder scheint die Sonne in Wald und Schlucht, See und Gletscher hinein, und an der neubeseelten Welt macht sich das helle Licht als ordnende, gruppirende, zeichnende Künstlerin zu schaffen; alles Nahe steigt ungeheuer im Werthe. „O Wunder! wo hatt’ ich die Augen?“ heisst es nun, und so leicht findet er des Staunens kein Ende, der selbst neubeseelte Mensch. Hat er doch immer noch das Nachbild jener älteren zwielichtfarbenen Welt in sich, von der aus alle Linien ins Unendliche hinüberzulaufen schienen und in die jeden Augenblick etwas Unendliches „hereinragte“, eine Welt in den Dämmerschatten des Imaginären und mit den Schlaglichtern des Jenseits, aber ohne die helle bestimmte Schönheit, deren allein die Realität en \Foreign{plein air} fähig ist. Ja, sie ist etwas gründlich Wunderbares, diese neugeschaffene, nein nur neuentdeckte \Sperr{empirische} Welt, die jetzt alles Menschliche ohne Ausnahme, auch das Reinste und Geistigste, in sich hineinzunehmen vermag — während jene andere „höhere“ Welt, in der man sich ehedem einen Theil der menschlichen Angelegenheiten verwaltet dachte, nun als chaotischer Abgrund zu unseren Füssen wogt. Werden wir diesem transcendenten Unhold, genannt „Etwas\Math{=}X“, noch weiterhin unser Bestes opfern? in diesen gurgelnden Schlund unsere Wünsche und Werthe, unsere goldenen Kleinode und Zierrathen hineinwerfen, in der phantastischen Hoffnung, dass er sie uns aufbewahre und dereinst vermehrt wiedergebe? Nein — wir, die am Rande des Abgrunds gehangen und uns schaudernd auf festen Boden zurückgerettet haben, wissen das Feste, Betretbare, den Erdgrund der Realität zu schätzen; uns, die den Grenzen der Menschheit nachgespürt und einen Fernblick auf \Foreign{ultima Thule} gethan haben, muss erlaubt sein, mit dem Jubel eines heimkehrenden Seefahrers den Rauch bewohnter Eilande aufsteigen zu sehen. Aber Jene, die nie über den Zaun ihres viereckigen Gemüsegärtchens hinausgekommen sind — dürfen sie ihre fröhliche Schollenkleberei und Sesshaftigkeit auch \Sperr{Optimismus} nennen? sollten sie, die nur ihre Winkelchen und ihre Winzigkeit und Specialität bejahen, unversehens unter die \Sperr{Weltbejahenden} gerathen sein?
\par

